% !TEX root = ../thesis.tex
% ==============================================================================
%  CHAPTER 2: PRELIMINARIES & FOUNDATIONS (فصل دوم: مفاهیم اولیه) (chapters/chapter2_preliminaries.tex)
% ==============================================================================
%  This chapter establishes the mathematical, theoretical, and algorithmic
%  foundations required for the thesis.
% ==============================================================================

\chapter{مفاهیم اولیه}\label{chap:chapter2_preliminaries}

این فصل به مرور مفاهیم ریاضی، مبانی نظری، معماری‌ها و ابزارهای پایه‌ای می‌پردازد که زیربنای ادبیات موضوع را تشکیل می‌دهند. هدف از ارائه‌ی این مباحث، پایه‌ریزی شالوده‌ی محاسباتی و نظری مورد نیاز برای درک صورت‌بندی مسئله در فصل~\ref{chap:chapter1_introduction} است، به‌گونه‌ای که خواننده‌ی آشنا با مبانی اولیه‌ی یادگیری ماشین و جبر خطی، پیش‌نیازهای لازم برای ورود به فصول آتی را در اختیار داشته باشد. لازم به ذکر است که بررسی تفصیلی پیشینه‌ی پژوهش، دسته‌بندی جامع روش‌های تفسیرپذیری و نقد مقایسه‌ای متدولوژی‌های موجود در فصل~\ref{chap:chapter3_literature} تشریح خواهد شد.

بخش~\ref{sec:2-1} به معرفی ساختار آماری مدل‌های زبانی بزرگ، فرایند تولید خودهمبسته، راهبردهای رمزگشایی و مقایسه‌ی پارادایم‌های جعبه‌سفید و جعبه‌سیاه اختصاص دارد. در بخش~\ref{sec:2-2}، معماری بنیادین ترنسفورمر، مکانیزم‌های توجه نقطه‌ای مقیاس‌یافته و چندسره، تفاوت‌های ساختاری خودتوجهی و توجه متقاطع، انواع ماسک‌گذاری (علّی و پدینگ) و ساختار لایه‌های باقیمانده تشریح می‌گردند.

سپس، در بخش~\ref{sec:2-3}، سطوح مختلف بازنمایی متن، از تعبیه‌های خرد توکنی و چالش‌های اختلال در زیرواژه‌ها تا رویکردهای مبتنی بر ارتقای واحد تحلیل به سطح جملات بر پایه‌ی اصول منطقی و شناختی، و در نهایت معماری شبکه‌های تعبیه‌ی جمله (\lr{Sentence-BERT}) بررسی می‌شوند. بخش~\ref{sec:2-4} به مدل‌سازی آماری از طریق مدل‌های مبتنی بر انرژی، چالش تحلیلی تابع افراز و روش‌های یادگیری تقابلی شامل تخمین تقابلی نویز (\lr{NCE}) و چارچوب \lr{InfoNCE} اختصاص دارد. در نهایت، در بخش~\ref{sec:2-5}، مبانی انتخاب ویژگی در سطح نمونه، صورت‌بندی ماسک‌های خلوت و ترفندهای آزادسازی پیوسته واکاوی خواهند شد.



\section{مدل‌های زبانی بزرگ}\label{sec:2-1}

مدل‌های زبانی بزرگ (\lr{LLM}) در سال‌های اخیر به استاندارد غالب در حوزه‌ی پردازش زبان طبیعی تبدیل شده‌اند. از منظر نظریه‌ی یادگیری آماری، این سیستم‌ها (نظیر خانواده‌ی مدل‌های \lr{GPT}~\cite{brown2020language} و \lr{LLaMA}~\cite{touvron2023llama}) توابع تخمین‌گر چگالی احتمال هستند که بر پایه‌ی پارادایم مدل‌سازی خودهمبسته عمل می‌کنند. در این دیدگاه، تولید متن به‌صورت یک فرایند تصادفی و گام‌به‌گام در طول زمان مدل‌سازی می‌شود.

فضای حالت پایه بر روی یک مجموعه‌ی متناهی از نمادهای گسسته به نام دایره‌ی واژگان تعریف می‌گردد که آن را با $\mathcal{V}$ نشان می‌دهیم ($|\mathcal{V}| < \infty$). اعضای این مجموعه، توکن‌ها (واحدهای پردازشی زبان اعم از کلمات کامل، زیرواژه‌ها یا نویسه‌ها) هستند. بستار کلینی بر روی این الفبا به‌صورت $\mathcal{V}^* = \bigcup_{t=0}^\infty \mathcal{V}^t$ تعریف می‌شود که نمایانگر مجموعه‌ی تمامی دنباله‌های متناهی از توکن‌هاست. 

فرض کنید دنباله‌ی توکن‌های پرامپت ورودی به طول $T_{\text{\lr{in}}}$ با $\mathbf{x} = (u_1, u_2, \dots, u_{T_{\text{\lr{in}}}}) \in \mathcal{V}^{T_{\text{\lr{in}}}} \subset \mathcal{V}^*$ و دنباله‌ی توکن‌های خروجی با $\mathbf{y} = (v_1, v_2, \dots, v_{T_{\text{\lr{out}}}}) \in \mathcal{V}^{T_{\text{\lr{out}}}} \subset \mathcal{V}^*$ مشخص شوند. یک مدل زبانی خودهمبسته، شبکه‌ی عصبی عمیقی با پارامترهای $\theta \in \Theta \subset \mathbb{R}^P$ است که توزیع احتمال شرطی توالی خروجی به شرط پرامپت ورودی، یعنی $P_\theta(\mathbf{y} \mid \mathbf{x})$، را تخمین می‌زند. با بهره‌گیری از قاعده‌ی ضرب احتمالات، این توزیع توأم به حاصل‌ضرب احتمالات شرطی یک‌گامی در طول زمان تجزیه می‌گردد~\cite{vaswani2017attention, brown2020language}:
\begin{equation}\label{eq:chap2_autoregressive}
    P_\theta(\mathbf{y} \mid \mathbf{x}) = \prod_{t=1}^{T_{\text{\lr{out}}}} P_\theta(v_t \mid \mathbf{x}, v_{<t})
\end{equation}
در این معادله، نماد $v_{<t} = (v_1, \dots, v_{t-1})$ بیانگر پیشوند توکن‌های تولیدشده در خروجی تا پیش از گام زمانی $t$ است، و برای گام اول، این پیشوند تهی در نظر گرفته می‌شود ($v_{<1} = \emptyset$). در هر گام زمانی $t$، شبکه، بافتار موجود شامل پرامپت ورودی و توکن‌های پیشین، یعنی $(\mathbf{x}, v_{<t})$، را پردازش کرده و بردار حالت پنهان آخرین لایه را به‌صورت $\mathbf{h}_t \in \mathbb{R}^{d_{\text{\lr{model}}}}$ استخراج می‌کند (که $d_{\text{\lr{model}}}$ ابعاد فضای پنهان مدل است). سپس از طریق ضرب این بردار در ماتریس وزن واژگانی $W_U \in \mathbb{R}^{|\mathcal{V}| \times d_{\text{\lr{model}}}}$، بردار امتیازات خام یا لوجیت‌ها به‌دست می‌آید:
\begin{equation}\label{eq:chap2_logits_generation}
    \mathbf{z}_t = W_U \mathbf{h}_t \in \mathbb{R}^{|\mathcal{V}|}
\end{equation}
که در آن، هر درایه‌ی $z_{t, v}$ نشان‌دهنده‌ی تمایل نسنجیده‌ی مدل به انتخاب توکن $v \in \mathcal{V}$ در گام زمانی $t$ است. برای تبدیل این امتیازات خام به یک توزیع احتمالاتی معتبر بر روی سیمپلکس دایره‌ی واژگان، عملگر سافت‌مکس بر روی این بردار اعمال می‌شود:
\begin{equation}\label{eq:chap2_softmax}
    P_\theta(v_t = v \mid \mathbf{x}, v_{<t}) = \text{\lr{Softmax}}{(\mathbf{z}_t)}_v = \frac{\exp(z_{t, v})}{\sum_{v' \in \mathcal{V}} \exp(z_{t, v'})} \quad \forall v \in \mathcal{V}
\end{equation}
در این رابطه، صورت کسر مقدار نمایی لوجیت توکن $v$ است و مخرج کسر حاصل‌جمع مقادیر نمایی تمامی توکن‌های دایره‌ی واژگان است تا تضمین شود مجموع احتمالات در سراسر واژگان دقیقاً برابر یک می‌گردد. پیامد ساختاری این معماری آن است که فضای جستجوی دنباله‌های خروجی به‌صورت نمایی با طول متن رشد می‌کند ($|\mathcal{V}|^{T_{\text{\lr{out}}}}$)؛ از این رو، ارزیابی کل‌نگر یک متن صرفاً از طریق ضرب احتمالات نقطه‌ای، به‌شدت مستعد انباشت خطاست.

فراتر از ابعاد محاسباتی، در سناریوهای کاربردی مدرن، فرایند تولید متن در مدل‌های زبانی از طریق یک راهبرد رمزگشایی مشخص انجام می‌پذیرد. برای مثال، در ساده‌ترین حالت که رمزگشایی حریصانه است، توکن خروجی در هر گام به‌صورت زیر برگزیده می‌شود:
\begin{equation}\label{eq:chap2_greedy}
    v_t = \arg\max_{v \in \mathcal{V}} P_\theta(v \mid \mathbf{x}, v_{<t})
\end{equation}
در رویکردهای تصادفی‌تر نظیر نمونه‌برداری هسته‌ای~\cite{holtzman2020curious} یا جستجوی پرتو، توکن از میان زیرمجموعه‌ای فیلترشده از محتمل‌ترین گزینه‌ها نمونه‌برداری یا جستجو می‌گردد. در تمامی این حالات، فرایند تولید تا جایی ادامه می‌یابد که مدل توکن توقف اختصاصی پایان توالی ($\langle\text{\lr{eos}}\rangle$) را صادر نماید.

از دیدگاه یک ناظر یا تحلیل‌گر بیرونی، این مکانیزم تولید باعث می‌شود که کل مدل زبانی بزرگ به‌صورت یک نگاشت اوراکل عمل کند:
\begin{equation}\label{eq:chap2_blackbox}
    f_\theta: \mathcal{V}^* \to \mathcal{V}^*, \quad \mathbf{y} = f_\theta(\mathbf{x})
\end{equation}
در این معادله، $\mathbf{x}$ پرامپت ورودی کاربر و $\mathbf{y}$ متن تولیدشده‌ی خروجی است. دسترسی به این نگاشت در ادبیات یادگیری ماشین در دو پارادایم اصلی قابل‌تفکیک است:
در رژیم جعبه‌سفید، تحلیل‌گر به تمامی اجزای گراف محاسباتی، ماتریس پارامترها ($\theta$)، لایه‌های پنهان، بردارهای فعال‌سازی ($\mathbf{h}_t$) و عملگرهای گرادیانی دسترسی کامل دارد. این دسترسی مستقیم به اطلاعات درونی، روش‌های جعبه‌سفید و مبتنی بر گرادیان را به خط‌پایه‌های مقایسه‌ای بنیادین و بسیار معتبری در تحلیل و ممیزی مدل‌ها بدل می‌سازد. در نقطه‌ی مقابل، در رژیم دسترسی جعبه‌سیاه (یا مدل پرس‌وجوی مرتبه‌صفر)، مدل در پسِ یک رابط برنامه‌نویسی (\lr{API}) مسدود قرار گرفته است و تحلیل‌گر به هیچ‌یک از سیگنال‌های درونی، وزن‌ها یا گرادیان‌ها دسترسی ندارد. در این شرایط، تنها عمل مجاز، ارسال پرامپت متنی گسسته‌ی $\mathbf{x}$ و دریافت پاسخ متنی محقق‌شده‌ی $\mathbf{y}$ است، بدون آنکه حتی بردار احتمالات خروجی (لوجیت‌ها) در دسترس قرار گیرد.

همان‌طور که اشنایدر (۲۰۲۴)~\cite{schneider2024explainable} در تبیین چشم‌انداز هوش مصنوعی مولد تفسیرپذیر تشریح می‌کند، ترکیب این انسداد ساختاری با ابعاد بالای فضای واژگان سبب می‌شود تا مسئله‌ی تفسیرپذیری، تبیین رفتار مدل و انتساب ویژگی‌ها در قبال یک پاسخ معین با چالش‌های نظری و محاسباتی عمیقی مواجه باشد. در حالی که روش‌های جعبه‌سفید با تکیه بر اطلاعات گرادیانی و فعال‌سازی‌های درونی به‌عنوان خط‌پایه‌های تحلیلی مهم در مدل‌های شفاف عمل می‌کنند، در سناریوهای جعبه‌سیاه عدم‌دسترسی به مشتقات و لوجیت‌ها، اجرای سرراست این رویکردها و همچنین روش‌های اختلالی سنتی را ناممکن می‌سازد. از این رو، واکاوی رفتار این سیستم‌ها نیازمند پایه‌ریزی مفاهیم ریاضی و معماری‌های جانشینی است که بتوانند صرفاً بر پایه‌ی تحلیل متن گسسته‌ی ورودی-خروجی، انتساب ویژگی‌ها را به‌شکلی ساختاریافته و معتبر ارزیابی نمایند.



\section{معماری ترنسفورمر}\label{sec:2-2}

معماری ترنسفورمر، که نخستین بار توسط واسوانی و همکاران (۲۰۱۷)~\cite{vaswani2017attention} معرفی گردید، به‌دلیل قابلیت پردازش موازی و ظرفیت بالا در مدل‌سازی وابستگی‌های بلندمدت، شالوده‌ی اصلی تمامی مدل‌های زبانی پیشرفته‌ی امروزی را تشکیل می‌دهد. در این بخش، اجزای بنیادین این معماری شامل روش‌های تزریق اطلاعات موقعیتی، مکانیزم‌های توجه نقطه‌ای و چندسره، انواع ماسک‌گذاری و بلوک‌های استاندارد ترنسفورمر به‌تفصیل فرمول‌بندی می‌شوند.


\subsection{رمزگذاری موقعیتی}\label{sec:2-2-1}

معماری ترنسفورمر برخلاف شبکه‌های عصبی بازگشتی\enfootnote{recurrent neural network} (\lr{RNN}) که متن را به‌صورت متوالی پردازش می‌کنند، فاقد ساختار ترتیبی درونی است و تمام توکن‌های ورودی را به‌طور هم‌زمان و موازی مورد محاسبه قرار می‌دهد. در نتیجه، مکانیزم توجه ذاتاً نسبت به‌جایگشت عناصر ورودی هم‌وردا\enfootnote{permutation equivariant} است؛ بدین‌معنا که تغییر در ترتیب کلمات یک جمله، بدون وجود اطلاعات جانبی، اثری بر محاسبات توجه نخواهد داشت.

برای جبران این کاستی و تزریق آگاهی ترتیبی به شبکه، اطلاعات جایگاه هر توکن در قالب رمزگذاری موقعیتی\enfootnote{positional encoding} به بردار بازنمایی اولیه‌ی آن اضافه می‌گردد. این رمزگذاری‌ها، بردارهایی با بعد دقیقاً برابر با بعد فضای پنهان مدل ($d_{\text{\lr{model}}}$) هستند که به‌صورت درایه‌به‌درایه با بردار تعبیه‌ی توکن‌ها جمع می‌شوند.

در فرمول‌بندی استاندارد تحلیلیِ ارائه‌شده توسط واسوانی و همکاران~\cite{vaswani2017attention}، از توابع موجی سینوسی و کسینوسی با بسامدهای متغیر در طول ابعاد استفاده می‌شود. در این رویکرد، برای یک موقعیت دلخواه در توالی $\mathit{pos} \in \{1, \dots, T\}$ و اندیس جفت‌بُعد $i \in \{0, 1, \dots, \lfloor d_{\text{\lr{model}}}/2 \rfloor - 1\}$، درایه‌های بردار موقعیت به‌صورت زیر محاسبه می‌گردند:

\begin{align}
    \text{\lr{PE}}_{(\mathit{pos}, 2i)} &= \sin\left(\frac{\mathit{pos}}{10000^{2i/d_{\text{\lr{model}}}}}\right) \label{eq:chap2_pe_sin} \\
    \text{\lr{PE}}_{(\mathit{pos}, 2i+1)} &= \cos\left(\frac{\mathit{pos}}{10000^{2i/d_{\text{\lr{model}}}}}\right) \label{eq:chap2_pe_cos}
\end{align}
در این معادلات، $\mathit{pos}$ نشان‌دهنده‌ی موقعیت مکانی توکن در توالی متنی، $i$ شاخص بعد در فضای فرکانس‌ها، و مقادیر مخرج یعنی $10000^{2i/d_{\text{\lr{model}}}}$ بیانگر طول موج‌هایی هستند که یک تصاعد هندسی را از $2\pi$ تا $10000 \cdot 2\pi$ تشکیل می‌دهند. این انتخاب تصادفی نبوده و از ویژگی‌های جبر مثلثاتی بهره می‌برد؛ بر اساس اتحادهای جمع زوایا، برای هر فاصله‌ی ثابت $k \in \mathbb{N}$، بردار موقعیت $\text{\lr{PE}}_{\mathit{pos}+k}$ می‌تواند به‌عنوان یک تبدیل خطی (دوران دوبعدی) از بردار $\text{\lr{PE}}_{\mathit{pos}}$ بازنمایی شود:\vspace{0.15cm}
\begin{equation}\label{eq:chap2_pe_rotation}
    \begin{pmatrix} 
        \text{\lr{PE}}_{(\mathit{pos}+k, 2i)} \\ 
        \text{\lr{PE}}_{(\mathit{pos}+k, 2i+1)} 
    \end{pmatrix} 
    = 
    \begin{pmatrix} 
        \cos(\omega_i k) & \sin(\omega_i k) \\ 
        -\sin(\omega_i k) & \cos(\omega_i k) 
    \end{pmatrix} 
    \begin{pmatrix} 
        \text{\lr{PE}}_{(\mathit{pos}, 2i)} \\ 
        \text{\lr{PE}}_{(\mathit{pos}, 2i+1)} 
    \end{pmatrix}
    \vspace{0.15cm}
\end{equation}
که در آن $\omega_i = 10000^{-2i/d_{\text{\lr{model}}}}$ بسامد زاویه‌ای است. این ویژگی هندسی به مدل اجازه می‌دهد تا بدون نیاز به یادگیری وزن‌های اضافی، فاصله‌ی نسبی میان عناصر را در کنار موقعیت مطلق آن‌ها به‌صورت خطی استخراج کند.

در کنار این رویکرد تحلیلی، روش متداول دیگر، استفاده از تعبیه‌های موقعیتی قابل‌یادگیری\enfootnote{learnable positional embeddings} است (که در مدل‌هایی نظیر \lr{BERT}~\cite{devlin2019bert} و \lr{GPT}~\cite{brown2020language} به‌کار می‌رود). در این رویکرد، برای هر جایگاه یک بردار پارامتریک در نظر گرفته می‌شود که به‌همراه سایر وزن‌های مدل در طول آموزش به‌روزرسانی می‌گردد. اگرچه تعبیه‌های قابل‌یادگیری انعطاف‌پذیری بالایی دارند، اما طول توالیِ قابل‌پردازش در آن‌ها به حداکثر طول دیده‌شده در فاز آموزش محدود است؛ در حالی که توابع سینوسی از منظر تئوری امکان تعمیم به طول‌های بیشتر را فراهم می‌سازند.


\subsection{مکانیزم توجه نقطه‌ای مقیاس‌یافته}\label{sec:2-2-2}

هسته‌ی پردازشی در معماری ترنسفورمر، مکانیزم توجه نقطه‌ای مقیاس‌یافته\enfootnote{scaled dot-product attention} است. این عملگر به‌عنوان یک ساختار بازیابی حافظه‌ی نرم و مشتق‌پذیر عمل می‌کند که در آن هر عضو از یک توالی از پرسمان‌ها، با مجموعه‌ای از کلیدها مقایسه شده و میانگین وزنیِ بردارهای ارزش متناظر محاسبه می‌گردد.

فرض کنید ماتریس پرسمان‌ها با $Q \in \mathbb{R}^{N_Q \times d_k}$، ماتریس کلیدها با $K \in \mathbb{R}^{N_K \times d_k}$ و ماتریس ارزش‌ها با $V \in \mathbb{R}^{N_K \times d_v}$ مشخص شوند (که در آن‌ها $N_Q$ طول توالی پرسمان، $N_K$ طول توالی منبع و $d_k, d_v$ ابعاد زیرفضاهای پرسمان و ارزش هستند). خروجی مکانیزم توجه به‌صورت زیر فرمول‌بندی می‌شود:
\begin{equation}\label{eq:chap2_attention}
    \text{\lr{Attention}}(Q, K, V) = \text{\lr{Softmax}}\left(\frac{QK^\top}{\sqrt{d_k}}\right)V
\end{equation}
در این معادله، حاصل‌ضرب ماتریسی $QK^\top \in \mathbb{R}^{N_Q \times N_K}$ ضرب داخلی میان تک‌تک پرسمان‌ها و کلیدها را محاسبه می‌کند که نمایانگر میزان شباهت و هم‌پوشانی آن‌هاست. ضریب مقیاس‌دهنده‌ی $\frac{1}{\sqrt{d_k}}$ دارای نقشی حیاتی در پایداری آموزش است: اگر مؤلفه‌های $Q$ و $K$ را متغیرهای تصادفی مستقل با میانگین صفر و واریانس یک فرض کنیم، ضرب داخلی آن‌ها واریانسی برابر با $d_k$ خواهد داشت. با بزرگ شدن ابعاد زیرفضا ($d_k$)، مقادیر حاصل از ضرب بسیار بزرگ شده و تابع سافت‌مکس را به نواحی اشباع سوق می‌دهند که در آن‌ها گرادیان‌ها بسیار ناچیز و نزدیک به صفر خواهند بود. تقسیم بر $\sqrt{d_k}$ واریانس را مجدداً به عدد یک بازمی‌گرداند و از پدیده‌ی محوشدگی گرادیان جلوگیری می‌کند~\cite{vaswani2017attention}. ماتریس حاصل از عملگر سافت‌مکس سطری، $A = \text{\lr{Softmax}}(QK^\top / \sqrt{d_k}) \in \mathbb{R}^{N_Q \times N_K}$، ماتریس وزن‌های توجه نامیده می‌شود که درایه‌های آن مقادیری در بازه‌ی $[0, 1]$ هستند و وزن اختصاص‌یافته به هر ارزش را مشخص می‌سازند.

برای ارتقای ظرفیت مدل در توجه هم‌زمان به جنبه‌های گوناگون اطلاعات در زیرفضاهای مختلف، این عملگر در قالب توجه چندسره\enfootnote{multi-head attention} (\lr{MHA}) گسترش می‌یابد. در این ساختار، ماتریس‌های ورودی با استفاده از ماتریس‌های تصویر پارامتریک به $h$ زیرفضای موازی تصویر می‌شوند:
\begin{align}
    \text{\lr{head}}_i &= \text{\lr{Attention}}(QW_i^Q, KW_i^K, VW_i^V) \quad \forall i \in \{1, \dots, h\} \label{eq:chap2_mha_head} \\
    \text{\lr{MHA}}(Q, K, V) &= \text{\lr{Concat}}(\text{\lr{head}}_1, \dots, \text{\lr{head}}_h)W^O \label{eq:chap2_mha_concat}
\end{align}
در این روابط، ماتریس‌های تبدیل خطی دارای ابعاد $W_i^Q \in \mathbb{R}^{d_{\text{\lr{model}}} \times d_k}$، $W_i^K \in \mathbb{R}^{d_{\text{\lr{model}}} \times d_k}$، و در نهایت $W_i^V \in \mathbb{R}^{d_{\text{\lr{model}}} \times d_v}$ و ماتریس خروجی دارای بعد $W^O \in \mathbb{R}^{hd_v \times d_{\text{\lr{model}}}}$ هستند. در طراحی‌های استاندارد، مقادیر ابعاد زیرفضاها به‌صورت $d_k = d_v = d_{\text{\lr{model}}} / h$ تنظیم می‌گردد تا بار محاسباتی کل این لایه‌ی چندسره با یک لایه‌ی توجه تک‌سره با بعد کامل $d_{\text{\lr{model}}}$ دقیقاً یکسان باقی بماند.


\subsection{خودتوجهی، توجه متقاطع و مکانیزم‌های ماسک‌گذاری}\label{sec:2-2-3}

بسته به خاستگاه داده‌های تشکیل‌دهنده‌ی پرسمان‌ها، کلیدها و ارزش‌ها، مکانیزم توجه به دو شکل کلی خودتوجهی\enfootnote{self-attention} و توجه متقاطع\enfootnote{cross-attention} پیکربندی می‌شود. همچنین برای کنترل جریان اطلاعات و جلوگیری از نقض قواعد ساختاری، از عملگرهای ماسک‌گذاری اختصاصی استفاده می‌گردد.

\paragraph{خودتوجهی.} در مکانیزم خودتوجهی، هر سه ماتریس $Q, K, V$ از یک توالی واحد (مانند ماتریس بازنمایی ورودی $X \in \mathbb{R}^{N \times d_{\text{\lr{model}}}}$) مشتق می‌شوند ($Q = XW^Q, K = XW^K, V = XW^V$). هدف از خودتوجهی، بازآرایی و غنی‌سازی بردار بازنمایی هر توکن بر پایه‌ی ارتباطات معنایی و نحوی آن با سایر توکن‌های همان توالی است.

نحوه‌ی اتصال نشانه‌ها در خودتوجهی به دو صورت است:
\begin{itemize}
    \item \textbf{خودتوجهی دوسویه\enfootnote{bidirectional self-attention}:} در مدل‌های رمزگذار (نظیر \lr{BERT}~\cite{devlin2019bert})، گراف توجه یک گراف کامل است و هر توکن مجاز است بدون محدودیت به تمامی توکن‌های پیشین و پسین در همان توالی توجه کند.
    \item \textbf{خودتوجهی علّی یا خودهمبسته\enfootnote{causal / autoregressive attention}:} در مدل‌های زایای خودهمبسته (نظیر \lr{GPT}~\cite{brown2020language})، مدل نباید در گام زمانی $t$ به اطلاعات توکن‌های آینده ($t+1, \dots$) دسترسی داشته باشد تا ساختار ترتیبی و پیش‌رونده‌ی تولید حفظ گردد. این قید ساختاری با افزودن یک ماتریس ماسک علّی به امتیازات توجه پیش از ورود به سافت‌مکس اعمال می‌شود.
\end{itemize}

\paragraph{توجه متقاطع.} در مکانیزم توجه متقاطع، پرسمان‌ها از یک توالی و کلیدها و ارزش‌ها از توالیِ دیگری تأمین می‌گردند. اگر $X \in \mathbb{R}^{N_X \times d_{\text{\lr{model}}}}$ توالی پرس‌وجوگر و $Y \in \mathbb{R}^{N_Y \times d_{\text{\lr{model}}}}$ توالی منبع باشد، ماتریس‌های توجه به فرم $Q = X W^Q$ و $K = Y W^K, V = Y W^V$ محاسبه می‌شوند. خروجی این لایه، توالی جدیدی به طول $N_X$ است که به ازای هر عنصر در $X$، اطلاعات هماهنگ و مرتبط را از توالی $Y$ تجمیع نموده است. در معماری‌های توالی‌به‌توالی سنتی، پرسمان‌ها از لایه‌های رمزگشا و کلیدها/ارزش‌ها از خروجی نهایی رمزگذار منشعب می‌شوند.

\paragraph{مکانیزم‌های ماسک‌گذاری در توجه.} در عمل، برای هدایت جریان اطلاعات و جلوگیری از اختلال در محاسبات، دو نوع ماسک‌گذاری اساسی در مکانیزم‌های توجه ترنسفورمر به‌کار گرفته می‌شوند:
\begin{enumerate}
    \item \textbf{ماسک علّی\enfootnote{causal mask}:} برای تضمین این اصل که احتمال تولید توکن بعدی صرفاً به پیشوند توکن‌های گذشته وابسته باشد، یک ماتریس ماسک مثلثی بالای $M_{\text{\lr{causal}}} \in \mathbb{R}^{N \times N}$ تعریف می‌شود:
    \begin{equation}\label{eq:chap2_causal_mask_def}
        {(M_{\text{\lr{causal}}})}_{ij} = \begin{cases} 0 & i \ge j \\ -\infty & i < j \end{cases}
    \end{equation}
    این ماتریس با حاصل‌ضرب داخلی مقیاس‌یافته جمع می‌گردد: $A = \text{\lr{Softmax}}\big(\frac{QK^\top}{\sqrt{d_k}} + M_{\text{\lr{causal}}}\big)$. از آنجا که $\exp(-\infty) = 0$ است، وزن توجه به تمامی توکن‌های آینده اکیداً صفر شده و هیچ‌گونه نشت اطلاعاتی از آینده رخ نمی‌دهد.
    \item \textbf{ماسک پدینگ\enfootnote{padding mask}:} در پردازش موازی دسته‌ای داده‌ها، طول جملات مختلف در یک دسته، یکسان نیست؛ بنابراین دنباله‌ها با افزودن توکن‌های بی‌اثر و ساختگی پدینگ ($\langle\text{\lr{pad}}\rangle$) به یک طول بیشینه‌ی مشترک رسانده می‌شوند. این توکن‌های ساختگی فاقد بار معنایی بوده و نباید در توزیع توجه مشارکت کنند. بدین منظور، یک ماتریس ماسک پدینگ $M_{\text{\lr{pad}}} \in {\{0, -\infty\}}^{N_Q \times N_K}$ تعریف می‌شود که درایه‌های متناظر با موقعیت‌های پدینگ در کلیدها ($K$) را برابر با $-\infty$ قرار می‌دهد:
    \begin{equation}\label{eq:chap2_pad_mask_def}
        {(M_{\text{\lr{pad}}})}_{ij} = \begin{cases} 0 & \text{اگر موقعیت } j \text{ در توالی منبع، یک توکن معتبر باشد} \\ -\infty & \text{اگر موقعیت } j \text{ در توالی منبع، یک توکن پدینگ باشد} \end{cases}
    \end{equation}
    اعمال این ماسک باعث می‌شود تا وزن توجه اختصاص‌یافته به توکن‌های پدینگ دقیقاً صفر گردد و این عناصر هیچ اثری بر روی بازنمایی نهایی توکن‌های معتبر بر جای نگذارند.
\end{enumerate}


\subsection{اتصالات باقیمانده و ساختار بلوک ترنسفورمر}\label{sec:2-2-4}

مکانیزم‌های توجه به‌تنهایی برای آموزش شبکه‌های متشکل از ده‌ها لایه‌ی متوالی کفایت نمی‌کنند؛ زیرا با عمیق‌تر شدن شبکه، انتشار سیگنال‌های گرادیان در فاز پس‌انتشار با افت و ناپایداری مواجه می‌شود. برای مهار این چالش، هر لایه‌ی ترنسفورمر مجهز به اتصالات باقیمانده\enfootnote{residual connections} است که ریشه در معماری \lr{ResNet}~\cite{he2016deep} دارد. در یک اتصال باقیمانده، بردار ورودی هر زیرلایه مستقیماً با بردار خروجی محاسبه‌شده توسط همان زیرلایه جمع می‌گردد: $\mathbf{x}_{\text{\lr{out}}} = \mathbf{x}_{\text{\lr{in}}} + \text{\lr{SubLayer}}(\mathbf{x}_{\text{\lr{in}}})$. این مسیر خطی و بدون مانع، یک شاهراه ارتباطی سرتاسری ایجاد می‌کند که در ادبیات تفسیرپذیری مکانیکی، «جریان باقیمانده»\enfootnote{residual stream} نامیده می‌شود~\cite{elhage2021mathematical}. این جریان به گرادیان‌ها اجازه می‌دهد تا بدون میرایی از عمیق‌ترین لایه‌ها به ابتدای شبکه بازگردند.

هر بلوک استاندارد در ترنسفورمر، علاوه بر لایه‌ی توجه، شامل یک شبکه‌ی عصبی پیش‌خور موقعیت‌محور\enfootnote{position-wise feed-forward network} (\lr{FFN}) است. این شبکه بر روی بردار بازنمایی تک‌تک توکن‌ها به‌صورت مستقل، موازی و با وزن‌های یکسان اعمال می‌شود. ساختار آن شامل دو تبدیل خطی متوالی با یک تابع فعال‌سازی غیرخطی در میان آن‌هاست. برای یک بردار فعال‌سازی $\mathbf{x} \in \mathbb{R}^{d_{\text{\lr{model}}}}$ داریم:
\begin{equation}\label{eq:chap2_ffn}
    \text{\lr{FFN}}(\mathbf{x}) = \max(0, \mathbf{x}W_1 + \mathbf{b}_1)W_2 + \mathbf{b}_2
\end{equation}
در این معادله، ماتریس وزن اول $W_1 \in \mathbb{R}^{d_{\text{\lr{model}}} \times d_{\mathit{ff}}}$ و بردار بایاس $\mathbf{b}_1 \in \mathbb{R}^{d_{\mathit{ff}}}$ بعد پنهان را موقتاً به ابعاد بزرگ‌تری منبسط می‌کنند (معمولاً $d_{\mathit{ff}} = 4d_{\text{\lr{model}}}$)، و سپس ماتریس وزن دوم $W_2 \in \mathbb{R}^{d_{\mathit{ff}} \times d_{\text{\lr{model}}}}$ و بایاس $\mathbf{b}_2 \in \mathbb{R}^{d_{\text{\lr{model}}}}$ آن را مجدداً به فضای پایه‌ی مدل بازمی‌گردانند. عملگر $\max(0, \cdot)$ نمایانگر تابع فعال‌سازی \lr{ReLU} است.

علاوه بر این، هر یک از زیرلایه‌های توجه و پیش‌خور، توسط عملیات نرمال‌سازی لایه‌ای\enfootnote{layer normalization} (\lr{LayerNorm})~\cite{ba2016layer} تثبیت می‌شوند. برخلاف نرمال‌سازی دسته‌ای که آمارها را در طول یک دسته محاسبه می‌کند، نرمال‌سازی لایه‌ای میانگین و واریانس را منحصراً در طول ابعاد ویژگیِ همان یک بردار فعال‌سازی استخراج می‌نماید. برای یک بردار فعال‌سازی $\mathbf{x} \in \mathbb{R}^d$ داریم:
\begin{equation}\label{eq:chap2_layernorm}
    \text{\lr{LN}}(\mathbf{x}) = \frac{\mathbf{x} - \mu}{\sqrt{\sigma^2 + \epsilon}} \odot \boldsymbol{\gamma} + \boldsymbol{\beta}
\end{equation}
که در آن $\mu = \frac{1}{d} \sum_{j=1}^d x_j$ میانگین درایه‌ها، $\sigma^2 = \frac{1}{d} \sum_{j=1}^d {(x_j - \mu)}^2$ واریانس درایه‌ها، $\boldsymbol{\gamma}, \boldsymbol{\beta} \in \mathbb{R}^d$ بردارهای مقیاس‌دهی و انتقالِ قابل‌یادگیری، و $\epsilon > 0$ یک عدد ثابت کوچک جهت پایداری عددی است. علامت $\odot$ نیز ضرب درایه‌به‌درایه‌ی هادامارد را نشان می‌دهد.

بدین‌ترتیب، جریان کامل محاسبات درون یک بلوک استاندارد رمزگذار ترنسفورمر برای ماتریس ورودی $X \in \mathbb{R}^{N \times d_{\text{\lr{model}}}}$ به‌صورت روابط متوالی زیر خلاصه می‌شود:
\begin{align}
    X^{\prime} &= \text{\lr{LN}}\big(X + \text{\lr{MHA}}(X, X, X)\big) \label{eq:chap2_transformer_block1} \\
    X^{\prime\prime} &= \text{\lr{LN}}\big(X^{\prime} + \text{\lr{FFN}}(X^{\prime})\big) \label{eq:chap2_transformer_block2}
\end{align}
این طراحی یکپارچه تضمین می‌کند که هر بلوک ترنسفورمر بتواند با استفاده از مکانیزم توجه، بافتار را از توکن‌های دیگر جمع‌آوری کرده و سپس با استفاده از لایه‌ی پیش‌خور، بازنمایی غیرخطی را به‌روزرسانی نماید.



\section{بازنمایی متنی و تعبیه‌ی جملات}\label{sec:2-3}

پردازش متون طبیعی توسط شبکه‌های عصبی عمیق، نیازمند نگاشت داده‌های نمادین و گسسته‌ی زبان به فضاهای برداری پیوسته است. در این بخش، ابتدا بازنمایی‌های توزیع‌شده در سطح کلمات بررسی شده و سپس چالش‌های نظری ناشی از اعمال اختلالات خرد در سطح توکن‌ها تشریح می‌گردد. در ادامه، مبانی منطقی و شناختی ارتقای واحد تحلیل به گزاره‌ها (جملات) واکاوی شده و در نهایت، معماری شبکه‌های تعبیه‌ی جمله و استراتژی‌های آموزش آن‌ها تبیین می‌شود.


\subsection{بازنمایی توزیع‌شده و نگاشت فضای گسسته}\label{sec:2-3-1}

در فرمول‌بندی نمادین اولیه، هر نشانه یا توکن $v$ از دایره‌ی واژگان $\mathcal{V}$ به‌صورت یک بردار یک‌گانه‌ی\enfootnote{one-hot vector} متعامد به نام $\mathbf{o}_v \in {\{0, 1\}}^{|\mathcal{V}|}$ نمایش داده می‌شود~\cite{manning2008introduction}. در این بردار، تنها درایه‌ی متناظر با اندیس آن نشانه برابر عدد یک و تمامی درایه‌های دیگر برابر صفر هستند. به‌دلیل ضرب داخلی صفر میان هر دو بردار یک‌گانه‌ی متمایز ($\mathbf{o}_i^\top \mathbf{o}_j = 0$ به ازای هر $i \neq j$)، این فضای برداری فاقد هرگونه مفهوم فاصله‌ی هندسی بوده و قادر به بازتاب شباهت‌های معنایی یا قرابت‌های کاربردی واژگان نیست.

برای حل این مسئله، از اصل «فرضیه‌ی توزیعی»\enfootnote{distributional hypothesis} و بازنمایی‌های توزیع‌شده\enfootnote{distributed representations}~\cite{mikolov2013efficient} بهره گرفته می‌شود. در این چارچوب، یک ماتریس تعبیه به نام $W_E \in \mathbb{R}^{d \times |\mathcal{V}|}$ تعریف می‌گردد که در آن $d \ll |\mathcal{V}|$ ابعاد فضای برداری پیوسته است. بردار پیوسته و متراکمِ متناظر با هر نشانه ($\mathbf{e}_v \in \mathbb{R}^d$) از طریق ضرب ماتریس تعبیه در بردار یک‌گانه استخراج می‌شود:
\begin{equation}\label{eq:chap2_embedding_lookup}
    \mathbf{e}_v = W_E \, \mathbf{o}_v
\end{equation}
در این هندسه‌ی برداری متراکم، مؤلفه‌های بردار موقعیت نشانه را در یک فضای پیوسته مشخص می‌کنند، به‌گونه‌ای که کلمات دارای کاربردهای مشابه یا بار معنایی نزدیک، در نقاطی با فاصله‌ی اقلیدسی اندک یا شباهت کسینوسی بالا نسبت به یکدیگر قرار می‌گیرند.


\subsection{توکن‌سازی زیرواژه‌ای و چالش خارج از توزیع}\label{sec:2-3-2}

مدل‌های زبانی مدرن برای کنترل حجم دایره‌ی واژگان و رفع مشکل مواجهه با کلمات خارج از دایره‌ی واژگان\enfootnote{out-of-vocabulary}، متن خام را به کلمات کامل تفکیک نمی‌کنند؛ بلکه از الگوریتم‌های توکن‌سازی زیرواژه‌ای نظیر رمزگذاری جفت‌بایت\enfootnote{byte-pair encoding} (\lr{BPE})~\cite{sennrich2016neural} بهره می‌برند. در این متدولوژی، کلمات به قطعات و پیشوندهای خردتر تجزیه می‌شوند تا مدل بتواند حتی کلمات جدید را بر پایه‌ی ترکیبی از زیرواژه‌ها بازشناسی کند.

اگرچه این رویکرد انعطاف‌پذیری محاسباتی بالایی به مدل می‌بخشد، اما در متدولوژی‌های تفسیرپذیریِ مبتنی بر اختلال یک نارسایی ساختاری جدی پدید می‌آورد. در روش‌های متداول ارزیابی پساعملی (نظیر چارچوب کلاسیک \lr{LIME}~\cite{ribeiro2016lime})، الگوریتم تلاش می‌کند تا با حذف، ماسک کردن یا جایگزینی تک‌تک توکن‌های ورودی، اهمیت هر یک را بر اساس تغییرات پدیدآمده در خروجی مدل بسنجد. فرض کنید یک پرامپت با حذف توکن $i$-ام مخدوش گردد:
\begin{equation}\label{eq:chap2_perturbation_basis}
    \mathbf{x}_{\setminus i} = (u_1, \dots, u_{i-1}, \text{\lr{[MASK]}}, u_{i+1}, \dots, u_{T_{\text{\lr{in}}}})
\end{equation}
از آنجا که بسیاری از توکن‌های حاصل از توکن‌سازی زیرواژه‌ای، واحدهای ناقص نحوی و املایی هستند (مانند سیلاب‌های منقطع از یک کلمه‌ی واحد)، حذف یا ماسک کردن آن‌ها ساختار گرامری جمله را متلاشی می‌کند. از منظر نظریه‌ی احتمال، توالی حاصل یک نمونه‌ی غیرطبیعی است که احتمال رخداد آن در توزیع داده‌های طبیعی زبان ($P_{\text{\lr{data}}}$) بسیار ناچیز بوده و به صفر میل می‌کند ($P_{\text{\lr{data}}}(\mathbf{x}_{\setminus i}) \to 0$)~\cite{feng2018pathologies}. 

پژوهش‌های فنگ و همکاران (۲۰۱۸)~\cite{feng2018pathologies} و مدسن و همکاران (۲۰۲۲)~\cite{madsen2022evaluating} نشان داده‌اند که این پدیده، مدل زبانی را با داده‌های \lr{OOD} مواجه می‌سازد. در چنین حالتی، نوسانات مشاهده‌شده در پاسخ مدل، ناشی از یک وابستگی معنایی واقعی میان محتوای ورودی و پاسخ مدل نیست؛ بلکه صرفاً واکنشی به اختلال نحوی و تخریب ساختار زبان است. این چالش نشان می‌دهد که ارزیابی انتساب ویژگی‌ها به‌صورت مستقیم در سطح زیرواژه‌ها می‌تواند به نتایجی کاذب و غیروفادارانه به فرایند تصمیم‌گیری مدل بینجامد.


\subsection{مبانی ارتقا به سطح جمله}\label{sec:2-3-3}

برای پرهیز از شکنندگی‌های توزیعی در سطح زیرواژه‌ها و دستیابی به واحدهایی با استقلال معنایی بالاتر، دسته‌ای از پژوهش‌های تحلیلی بر این نکته تأکید دارند که واحد تحلیل به‌جای نشانه‌های خرد نحوی، به مفاهیم معنایی مستقل نظیر گزاره‌ها یا جملات کامل ارتقا یابد. در این ساختار، هر جمله به‌عنوان یک گزاره‌ی زبانی کامل، نقشی متمایز را به‌عنوان یک پیش‌آیند بالقوه در بافتار پرامپت ایفا می‌کند. این رویکرد بر پایه‌های نظری شناخته‌شده‌ای در فلسفه‌ی تحلیلی زبان و روان‌شناسی شناختی استوار است:

در فلسفه‌ی زبان، گوتلوب فرگه در کتاب بنیادین «مبانی حساب» (۱۸۸۴)~\cite{frege1991foundations}، «اصل بافتار» را صورت‌بندی نمود. بر اساس این اصل، کلمات منفرد در انزوا فاقد معنای مستقل و متعین هستند و هویت معنایی واقعی آن‌ها تنها در ساختار یک گزاره‌ی کامل شکل می‌گیرد. در زبان‌شناسی محاسباتی، این ایده را می‌توان بدین‌صورت تبیین کرد که اگر $\mathcal{V}$ فضای نشانه‌ها و $\mathcal{S} \subset \mathcal{V}^*$ فضای جملات معتبر باشد، یک توکن منفرد $v \in \mathcal{V}$ به‌تنهایی نمی‌تواند موضوع یک گزاره‌ی صدق قرار گیرد؛ بدین‌معنا که نمی‌توان یک تابع ارزش صدق $\nu: \mathcal{V} \to \{0, 1\}$ برای کلمه‌ای مجزا تعریف کرد. در مقابل، جمله ($S \in \mathcal{S}$) کوچک‌ترین واحد حامل یک اندیشه و گزاره‌ی کامل است که قابلیت پذیرش ارزش صدق منطقی را داراست ($\nu(S) \in \{0, 1\}$). از این دیدگاه، معنای یک نشانه‌ی منفرد ($\mathcal{M}(v)$) مفهومی چندوجهی است که صرفاً با شرطی‌سازی بر بافتار یک گزاره‌ی کامل تثبیت می‌گردد: $\mathcal{M}(v \mid S)$.

این ساختار گزاره‌ای در مدل‌های پردازش شناختی نیز تأیید شده است. والتر کینتش (۱۹۹۸)~\cite{kintsch1998comprehension} در مدل شناختی «ساخت‌وساز-ادغام»، تبیین می‌کند که ذهن انسان در هنگام ادراک مطلب، متن را به‌صورت توالی صرف از کلمات در حافظه نگهداری نمی‌کند. در ذهن، یک شبکه‌ی گزاره‌ای تحت عنوان «پایگاه متن»\enfootnote{textbase} به‌شکل یک گراف معنایی $G = (V_G, E_G)$ تشکیل می‌شود. در این گراف، گره‌ها ($V_G \subset \mathcal{S}$) گزاره‌های استخراج‌شده از جملات هستند و یال‌ها ($E_G$) روابط منطقی، ارجاعی و سببی میان این ایده‌ها را تشکیل می‌دهند.

بر پایه‌ی این شواهد، استفاده از جمله به‌عنوان واحد بررسی، این مزیت را دارد که با حذف یا ارزیابی یک واحد، ساختار گرامری درونی آن آسیب ندیده و شبکه با یک گزاره‌ی کامل، به‌جای یک قطعه‌ی نحوی ناقص، مواجه می‌شود. بدین‌ترتیب، جملات به‌عنوان پیش‌آیندهای ورودی، بازنمایی ساختاریافته‌تر و انسان‌خواناتری برای تحلیل انتساب و ردیابی تأثیر فراهم می‌سازند.


\subsection{معماری شبکه‌های تعبیه‌ی جمله}\label{sec:2-3-4}

به‌کارگیری جملات به‌عنوان ورودی‌های پردازش عصبی، نیازمند تابعی است که بتواند یک جمله با طول متغیر از توکن‌ها $S = (u_1, u_2, \dots, u_k) \in \mathcal{V}^k$ را به یک بردار متراکم با طول ثابت $\mathbf{v}_S \in \mathbb{R}^d$ نگاشت کند ($f_{\text{\lr{sent}}}: \mathcal{S} \to \mathbb{R}^d$). چارچوب \lr{Sentence-BERT}~\cite{reimers2019sentencebert} با بهره‌گیری از شبکه‌های عصبی سیامی\enfootnote{Siamese networks} این نگاشت را فرمول‌بندی نموده است.

در این معماری، توکن‌های جمله ابتدا از یک ترنسفورمر پایه (نظیر \lr{BERT}~\cite{devlin2019bert} یا \lr{RoBERTa}~\cite{liu2021roberta}) عبور داده می‌شوند. خروجی، دنباله‌ای از بردارها به‌صورت ماتریس $\mathbf{H} = (\mathbf{h}_1, \dots, \mathbf{h}_k) \in \mathbb{R}^{k \times d}$ است که در آن هر $\mathbf{h}_i$ بازنمایی بافتاریافته‌ی توکن $i$-ام را نشان می‌دهد. برای فشرده‌سازی این ماتریس با ابعاد متغیر به یک بردار متراکم واحد، از یک لایه‌ی تجمیع استفاده می‌گردد.

در ادبیات این حوزه، راهبردهای متعددی برای این تبدیل بررسی شده‌اند، از جمله استفاده از بازنمایی توکن آغازین ($\text{\lr{[CLS]}}$)، تجمیع بیشینه\enfootnote{max pooling} و تجمیع میانگین\enfootnote{mean pooling}. ارزیابی‌های تجربی ریمرز و گورویچ نشان دادند که تجمیع میانگین در طول بعد توالی، باثبات‌ترین و دقیق‌ترین بازنمایی را تولید می‌کند. بردار حاصل از این تجمیع به‌صورت زیر به‌دست می‌آید:
\begin{equation}\label{eq:chap2_mean_pooling}
    \mathbf{v}_S = \frac{1}{k} \sum_{i=1}^k \mathbf{h}_i
\end{equation}
در این معادله، $k$ تعداد توکن‌های جمله‌ی $S$ و $\mathbf{h}_i$ بردار خروجی توکن $i$-ام از ترنسفورمر پایه است. حاصل این عملیات، یک بردار واحد $\mathbf{v}_S \in \mathbb{R}^d$ است که میانگین بار معنایی توکن‌های جمله را در خود جای داده است.

برای اینکه فضای برداری حاصل ساختار هندسی معناداری پیدا کند، شبکه‌ی \lr{Sentence-BERT} بر روی مجموعه‌داده‌های مقیاس‌بزرگ با دو هدف آموزشی کلیدی تنظیم ظریف می‌شود: در وظیفه‌ی نخست، یعنی استنتاج زبان طبیعی\enfootnote{natural language inference} (\lr{NLI})~\cite{bowman2015large, williams2018broad}، مدل باید نوع رابطه‌ی منطقی میان دو جمله (شامل استلزام، تناقض یا خنثی) را تشخیص دهد. برای این منظور، بردارهای دو جمله یعنی $\mathbf{u}$ و $\mathbf{v}$ استخراج شده و به‌همراه بردار تفاضل مطلق آن‌ها، یک بردار الحاقی به‌صورت $(\mathbf{u}, \mathbf{v}, |\mathbf{u} - \mathbf{v}|) \in \mathbb{R}^{3d}$ تشکیل می‌گردد. این بردار سپس در یک ماتریس وزن ضرب شده و با تابع زیان آنتروپی متقاطع آموزش می‌بیند تا شبکه تفاوت‌های ظریف گزاره‌ای را یاد بگیرد. در وظیفه‌ی دوم، یعنی شباهت معنایی متون\enfootnote{semantic textual similarity} (\lr{STS})~\cite{cer2017semeval}، هدف، تطبیق هندسی بردارها با قضاوت انسان درباره‌ی شباهت جملات است. در اینجا، شباهت کسینوسی میان دو بردار جمله محاسبه شده و با استفاده از تابع زیان میانگین مربعات خطا، به امتیازات برچسب‌گذاری‌شده نزدیک می‌گردد:
\begin{equation}\label{eq:chap2_cosine_sim}
    S_C(\mathbf{v}_{S_A}, \mathbf{v}_{S_B}) = \frac{\mathbf{v}_{S_A}^\top \mathbf{v}_{S_B}}{\|\mathbf{v}_{S_A}\|_2 \|\mathbf{v}_{S_B}\|_2}
\end{equation}
در این فرمول، صورت کسر ضرب داخلی میان دو بردار جمله است و مخرج کسر حاصل‌ضرب نُرم اقلیدسی آن‌هاست. این سنجه زاویه‌ی میان دو بردار را در فضای چندبعدی اندازه‌گیری کرده و مقداری بین $-1$ و $+1$ تولید می‌کند که مستقل از طول مطلق بردارهاست.

این فرایند آموزش سبب می‌شود تا ویژگی‌های هندسی تعبیه‌ها هم‌راستا با درک انسانی سازمان‌دهی شوند. در همین راستا، خطیر و همکاران (۲۰۲۵)~\cite{khatir2025aligned} با تحلیل‌های مبتنی بر گراف فازی و نظریه‌ی جوامع نشان دادند که این‌گونه تعبیه‌های اولیه، دارای خاصیتی تحت عنوان «هم‌تراز در مبدأ»\enfootnote{aligned at the start} هستند؛ بدین‌معنا که خوشه‌بندی‌های معنایی و ساختارهای مفهومی به‌طور ارگانیک حتی پیش از اعمال تنظیمات خاص، در این فضا شکل گرفته‌اند.



\section{مدل‌های مبتنی بر انرژی}\label{sec:2-4}

مدل‌های مبتنی بر انرژی (\lr{EBM}) یک چارچوب ریاضی عام و انعطاف‌پذیر برای یادگیری ماشین ارائه می‌دهند که در آن وابستگی و انطباق‌پذیری میان متغیرها از طریق نسبت‌دادن یک کمیت اسکالر نرمال‌نشده به نام «انرژی» ارزیابی می‌گردد~\cite{lecun2006tutorial}. برخلاف مدل‌های خودهمبسته که متغیرها را به‌صورت گام‌به‌گام و مقید به نرمال‌سازی محلی پیش‌بینی می‌کنند، یک مدل انرژی می‌تواند ورودی و خروجی را به‌صورت ساختاری و یکپارچه داوری کند. در این بخش، مبانی نظری، چالش‌های تحلیلی آموزش و متدولوژی‌های بهینه‌سازی تقابلی در این مدل‌ها تشریح می‌شوند.


\subsection{فرمول‌بندی پایه و توزیع گیبس}\label{sec:2-4-1}

از منظر صوری، یک مدل مبتنی بر انرژی تابعی اسکالر و مشتق‌پذیر به فرم زیر است:
\begin{equation}\label{eq:chap2_ebm_func_def}
    \mathcal{E}_\theta : \mathcal{X} \times \mathcal{Y} \to \mathbb{R}
\end{equation}
که در آن $\mathcal{X}$ فضای متغیرهای ورودی، $\mathcal{Y}$ فضای متغیرهای خروجی، و $\theta \in \mathbb{R}^P$ بردار پارامترهای مدل است. تابع $\mathcal{E}_\theta(x, y)$ مقیاسی از عدم‌انطباق یا ناهماهنگی میان پیکربندی متغیر ورودی $x$ و متغیر خروجی $y$ را نشان می‌دهد. بر این اساس، مقادیر پایین‌تر انرژی نمایانگر انطباق‌پذیری بیشتر و مطلوبیت بالاتر آن پیکربندی است، و در مقابل، مقادیر بالاتر انرژی نشان‌دهنده‌ی تناقض یا عدم‌ارتباط میان متغیرهاست.

برای پیوند دادن این اسکالر نرمال‌نشده با نظریه‌ی احتمال، از ساختار کلاسیک مکانیک آماری و توزیع گیبس-بولتزمن\enfootnote{Gibbs-Boltzmann distribution} بهره گرفته می‌شود. توزیع احتمال شرطی خروجی $y \in \mathcal{Y}$ به شرط ورودی $x \in \mathcal{X}$ به‌صورت زیر تعریف می‌گردد:
\begin{equation}\label{eq:chap2_ebm_prob}
    P_\theta(y \mid x) = \frac{\exp\big( -\mathcal{E}_\theta(x, y) / \tau \big)}{Z_\theta(x)}
\end{equation}
در این رابطه، متغیر $\tau > 0$ پارامتر دما است که میزان تمرکز یا تیزی توزیع را کنترل می‌کند؛ با میل کردن $\tau \to 0$، توزیع بر روی کمینه‌ی مطلق انرژی تمرکز می‌یابد و با افزایش $\tau \to \infty$، توزیع به سمت یک توزیع یکنواخت میل می‌کند. مخرج کسر، یعنی $Z_\theta(x)$، «تابع افراز»\enfootnote{partition function} نامیده می‌شود که عمل نرمال‌سازی را انجام می‌دهد:
\begin{equation}\label{eq:chap2_partition_func}
    Z_\theta(x) = \sum_{y' \in \mathcal{Y}} \exp\big( -\mathcal{E}_\theta(x, y') / \tau \big) \quad \text{یا} \quad \int_{\mathcal{Y}} \exp\big( -\mathcal{E}_\theta(x, y') / \tau \big) \, dy'
\end{equation}
که در آن حاصل‌جمع (یا انتگرال) بر روی تمامی پیکربندی‌های ممکن در فضای خروجی $\mathcal{Y}$ محاسبه می‌شود تا تضمین گردد مجموع احتمالات برابر با یک است.


\subsection{بن‌بست محاسباتی تابع افراز}\label{sec:2-4-2}

در دامنه‌های پردازش زبان طبیعی و تولید متن، خروجی $y$ یک دنباله‌ی متنی متشکل از $T$ توکن است که هر یک از دایره‌ی واژگان $\mathcal{V}$ انتخاب می‌شوند. در نتیجه، فضای حالات خروجی دارای بعد نمایی است:
\begin{equation}\label{eq:chap2_state_space_cardinality}
    |\mathcal{Y}| = |\mathcal{V}|^T
\end{equation}
محاسبه‌ی صریح تابع افراز $Z_\theta(x)$ مستلزم محاسبه‌ی این جمع نمایی بر روی تمامی ترکیبات ممکن است که از نظر محاسباتی کاملاً حل‌ناپذیر است.

این بن‌بست تحلیلی زمانی بحرانی می‌شود که بخواهیم مدل را با روش‌های سنتی مبتنی بر درست‌نمایی بیشینه آموزش دهیم. در این استراتژی، هدف کمینه‌سازی منفیِ لگاریتم درست‌نمایی\enfootnote{negative log-likelihood} (\lr{NLL}) بر روی توزیع داده‌های واقعی ($P_{\text{\lr{data}}}$) است:
\begin{equation}\label{eq:chap2_nll_loss}
    \mathcal{L}_{\text{\lr{NLL}}}(\theta) = -\mathbb{E}_{(x, y) \sim P_{\text{\lr{data}}}} \big[ \log P_\theta(y \mid x) \big]
\end{equation}
با جایگذاری رابطه‌ی توزیع گیبس (معادله‌ی~\ref{eq:chap2_ebm_prob})، عبارت درون امید ریاضی به‌صورت زیر بسط داده می‌شود:
\begin{equation}\label{eq:chap2_log_prob_expansion}
    \log P_\theta(y \mid x) = -\frac{1}{\tau} \mathcal{E}_\theta(x, y) - \log Z_\theta(x)
\end{equation}
برای محاسبه‌ی گرادیان این تابع زیان نسبت به پارامترهای مدل ($\theta$)، ابتدا مشتق عبارت دوم را بررسی می‌کنیم:
\begin{align}
    \nabla_\theta \log Z_\theta(x) &= \frac{\nabla_\theta Z_\theta(x)}{Z_\theta(x)} \nonumber \\
    &= \frac{1}{Z_\theta(x)} \sum_{y' \in \mathcal{Y}} \nabla_\theta \exp\big( -\mathcal{E}_\theta(x, y') / \tau \big) \nonumber \\
    &= \frac{1}{Z_\theta(x)} \sum_{y' \in \mathcal{Y}} \exp\big( -\mathcal{E}_\theta(x, y') / \tau \big) \left( -\frac{1}{\tau} \nabla_\theta \mathcal{E}_\theta(x, y') \right) \nonumber \\
    &= -\frac{1}{\tau} \sum_{y' \in \mathcal{Y}} P_\theta(y' \mid x) \nabla_\theta \mathcal{E}_\theta(x, y') \nonumber \\
    &= -\frac{1}{\tau} \mathbb{E}_{y' \sim P_\theta(\cdot \mid x)} \big[ \nabla_\theta \mathcal{E}_\theta(x, y') \big] \label{eq:chap2_grad_partition_derivation}
\end{align}

در این اشتقاق، عبارت $\frac{\exp(-\mathcal{E}_\theta(x, y') / \tau)}{Z_\theta(x)}$ دقیقاً همان احتمال $P_\theta(y' \mid x)$ است که جمع وزنی را به یک امید ریاضی بر روی توزیع مدل تبدیل می‌کند. با اعمال این نتیجه در مشتق معادله‌ی~\ref{eq:chap2_nll_loss}، گرادیان نهایی زیان درست‌نمایی به فرم مشهور زیر حاصل می‌شود:
\begin{equation}\label{eq:chap2_mle_gradient}
    \nabla_\theta \mathcal{L}_{\text{\lr{NLL}}}(\theta) = \frac{1}{\tau} \bigg( \underbrace{\mathbb{E}_{(x, y) \sim P_{\text{\lr{data}}}} \big[\nabla_\theta \mathcal{E}_\theta(x, y)\big]}_{\text{فاز مثبت (کاهش انرژی نمونه‌های داده)}} - \underbrace{\mathbb{E}_{x \sim P_{\text{\lr{data}}}} \mathbb{E}_{y' \sim P_\theta(\cdot \mid x)} \big[\nabla_\theta \mathcal{E}_\theta(x, y')\big]}_{\text{فاز منفی (افزایش انرژی نمونه‌های مدل)}} \bigg)
\end{equation}

این معادله نشان می‌دهد که بهینه‌سازی درست‌نمایی از دو فاز متضاد تشکیل یافته است:
\begin{itemize}
    \item \textbf{فاز مثبت:} گرادیان روی داده‌های واقعی محاسبه می‌شود و هدف آن کاهش سطح انرژی برای زوج‌های منطبق مشاهده‌شده در مجموعه‌داده است. این عبارت مستقیماً بر روی نمونه‌های آموزشی قابل‌محاسبه است.
    \item \textbf{فاز منفی:} گرادیان روی نمونه‌هایی محاسبه می‌شود که توسط خود مدل ($P_\theta$) تولید می‌شوند و هدف آن بالا بردن سطح انرژی در سایر نقاط فضا است تا مدل به پاسخ‌های نامعتبر انرژی کمی اختصاص ندهد.
\end{itemize}
چالش اصلی دقیقاً در همین فاز منفی نهفته است؛ محاسبه‌ی این امید ریاضی مستلزم نمونه‌برداری از توزیع نرمال‌نشده‌ی مدل ($P_\theta$) است. در فضاهای پیوسته، تقریب این توزیع از طریق روش‌های مبتنی بر گرادیان نظیر دینامیک لانژون انجام می‌شود، اما در فضای گسسته‌ی متن، نمونه‌برداری با روش‌های زنجیره‌ی مارکوف مونت‌کارلو (\lr{MCMC}) نیازمند ارزیابی تمام کلمات واژگان در هر موقعیت از توالی است. این امر سربار محاسباتی را به‌شکل فلج‌کننده‌ای افزایش می‌دهد~\cite{lecun2006tutorial, deng2020residual, qin2022cold}.

با وجود این پیچیدگی در آموزش، یک مزیت بنیادین در مدل‌های انرژی وجود دارد: \textbf{تابع افراز $Z_\theta(x)$ صرفاً برای نرمال‌سازی احتمالات در فاز آموزش مورد نیاز است اما بنا به تعریف این مدل‌ها، تابع در فاز استنتاج و ارزیابی به‌طور کامل کنار می‌رود.} 

در زمان استنتاج، پیش‌بینی بهینه‌ترین خروجی یا ارزیابی میزان انطباق‌پذیری یک زوج داده‌شده، مستقیماً از طریق کمینه‌سازی انرژی یا محاسبه‌ی یک گذر پیش‌رو ساده انجام می‌شود:
\begin{equation}\label{eq:chap2_ebm_inference}
    y^* = \arg\min_{y \in \mathcal{Y}} \mathcal{E}_\theta(x, y)
\end{equation}
در این معادله، $y^*$ بیان‌گر خروجی‌ای است که کمینه‌ترین انرژی را تولید می‌کند. این ویژگی به مدل اجازه می‌دهد تا بدون نیاز به نرمال‌سازی احتمالات بر روی تک‌تک توکن‌ها، داده را به‌عنوان یک ساختار کلان ارزیابی کند تا نواحی کم‌انرژی آن را شناسایی کند. از این رو، جامعه‌ی پژوهشی به سمت استراتژی‌های آموزشیِ بدون نیاز به محاسبه‌ی تابع افراز حرکت کرده است.


\subsection{شبکه‌های انرژی ساختاریافته}\label{sec:2-4-3}

برای سال‌ها، مدل‌سازی وابستگی‌ها در خروجی‌های ساختاریافته توسط میدان‌های تصادفی شرطی\enfootnote{conditional random fields} (\lr{CRF})~\cite{lafferty2001conditional} انجام می‌گرفت. در یک \lr{CRF}، برای اینکه محاسبه‌ی تابع افراز از طریق الگوریتم‌های برنامه‌ریزی پویا (نظیر الگوریتم پیشرو-پسرو) امکان‌پذیر باشد، تابع انرژی کل به مجموعه‌ای از پتانسیل‌های محلی روی کلیک‌های\enfootnote{cliques} کوچک گراف تجزیه می‌شد:
\begin{equation}\label{eq:chap2_crf_energy}
    \mathcal{E}_{\text{\lr{CRF}}}(x, y) = \sum_{c \in \mathcal{C}} \mathcal{E}_c(x, y_c)
\end{equation}
در این معادله، $\mathcal{C}$ مجموعه‌ی کلیک‌های گراف، $y_c$ زیرمجموعه‌ی متغیرهای خروجی درون هر کلیک، و $\mathcal{E}_c$ تابع پتانسیل محلی آن کلیک است. اگرچه این تجزیه‌ی خطی مدل را قابل‌حل می‌ساخت، اما فرض استقلال شرطی مارکوف باعث می‌شد که مدل نتواند تعاملات پیچیده، غیرخطی و دوربرد موجود در زبان طبیعی را پوشش دهد.

برای عبور از این محدودیت، بلانجر و مک‌کالوم (۲۰۱۶)~\cite{belanger2016spen} چارچوب «شبکه‌های انرژی پیش‌بینی ساختاریافته»\enfootnote{structured prediction energy networks} (\lr{SPEN}) را معرفی نمودند. در این معماری، فرضیات خطی و تجزیه‌های مارکوف کنار گذاشته شده و تابع انرژی سراسری به‌صورت مستقیم توسط یک شبکه‌ی عصبی عمیق مدل‌سازی می‌شود:
\begin{equation}\label{eq:chap2_spen_energy}
    \mathcal{E}_{\text{\lr{SPEN}}}(x, y) = \text{\lr{DeepNN}}_\theta\big(\mathbf{f}_x(x), \mathbf{f}_y(y)\big)
\end{equation}
در این فرمول‌بندی، $\mathbf{f}_x(x) \in \mathbb{R}^d$ و $\mathbf{f}_y(y) \in \mathbb{R}^d$ توابع نگاشت پیوسته‌ای هستند که ویژگی‌های ورودی و ساختار کامل خروجی را استخراج می‌کنند، و $\text{\lr{DeepNN}}_\theta$ یک شبکه‌ی عصبی عمیق با پارامترهای $\theta$ است که این دو بردار را به‌عنوان ورودی دریافت کرده و انطباق‌پذیری کلان آن‌ها را در قالب یک عدد اسکالر واحد داوری می‌کند. این چارچوب اثبات کرد که می‌توان ساختارهای خروجی پیچیده را فراتر از پتانسیل‌های مجزا، به‌صورت کل‌نگر و یکپارچه ارزیابی کرد.


\subsection{یادگیری تقابلی و بهینه‌سازی انرژی}\label{sec:2-4-4}

برای حل بن‌بست فاز منفی در معادله‌ی~\ref{eq:chap2_mle_gradient} و اجتناب از محاسبه‌ی تابع افراز، پارادایم یادگیری تقابلی\enfootnote{contrastive learning} به‌عنوان یک راهکار بهینه‌سازی کارآمد مطرح می‌گردد~\cite{lecun2006tutorial}.

هدف اصلی در یادگیری تقابلی، شکل‌دهی هندسی به چشم‌انداز انرژی بدون نیاز به مدل‌سازی توزیع روی کل فضای $\mathcal{Y}$ است. در این استراتژی، به‌جای محاسبه‌ی انتگرال روی تمامی حالات نامتناهی، مدل در حین آموزش به‌صورت هم‌زمان با دو دسته از داده‌ها تغذیه می‌شود: نمونه‌های مثبت $y^+$ که نشان‌دهنده‌ی خروجی‌های واقعی و منطبق بر ورودی $x$ هستند، و نمونه‌های منفی $y^-$ که از یک توزیع نویز یا اختلال ($P_n$) نمونه‌برداری می‌شوند و نمایانگر حالت‌های نامعتبر یا ناسازگارند.

یکی از متداول‌ترین روش‌های پیاده‌سازی این ایده، استفاده از تابع زیان حاشیه‌محور\enfootnote{margin loss / hinge loss} است که ریشه در پژوهش‌های هدسل و همکاران (۲۰۰۶)~\cite{hadsell2006dimensionality} و لکان و همکاران (۲۰۰۶)~\cite{lecun2006tutorial} دارد:
\begin{equation}\label{eq:chap2_margin_loss}
    \mathcal{L}_{\text{\lr{margin}}}(\theta) = \mathbb{E}_{(x, y^+) \sim P_{\text{\lr{data}}}, \, y^- \sim P_n(\cdot \mid x)} \Big[ \max\big(0, \, \mathcal{E}_\theta(x, y^+) - \mathcal{E}_\theta(x, y^-) + \gamma\big) \Big]
\end{equation}
در این رابطه، $\mathcal{E}_\theta(x, y^+)$ انرژی زوج مثبت، $\mathcal{E}_\theta(x, y^-)$ انرژی زوج منفی، و $\gamma > 0$ حاشیه‌ی امن مطلوب است. عملگر $\max(0, \cdot)$ تضمین می‌کند که اگر انرژی نمونه‌ی منفی حداقل به اندازه‌ی $\gamma$ واحد از انرژی نمونه‌ی مثبت بزرگ‌تر باشد ($\mathcal{E}_\theta(x, y^-) \ge \mathcal{E}_\theta(x, y^+) + \gamma$)، مقدار زیان صفر شود. این تابع زیان، یک شیب گرادیانی ایجاد می‌کند که خمینه‌ی داده‌های معتبر را به سمت انرژی‌های پایین کشیده و نمونه‌های مخدوش را به سطوح بالای انرژی می‌راند. با این وجود، توابع حاشیه‌ای ساده به‌دلیل حساسیت بالا به انتخاب ابرپارامتر $\gamma$ و ناپایداری در فضاهای پیچیده، در پردازش زبان طبیعی غالباً جای خود را به روش‌های احتمالاتی تقابلی می‌دهند.


\subsection{تخمین نویز تقابلی و نقش گرادیان‌ها}\label{sec:2-4-5}

یک رویکرد بسیار پایدارتر و اصولی‌تر برای آموزش مدل‌های مبتنی بر انرژی، بهره‌گیری از روش تخمین تقابلی نویز\enfootnote{noise-contrastive estimation} (\lr{NCE}) است که توسط گوتمان و هیواریانن (۲۰۱۰)~\cite{gutmann2010noise} جهت تخمین پارامترهای مدل‌های آماریِ فاقد ضریب نرمال‌سازی مطرح شد.

ایده‌ی پایه‌ای \lr{NCE} آن است که به‌جای تلاش برای تقریب زدن انتگرال ناممکن در مخرج کسر تابع افراز، مسئله‌ی تخمین چگالی احتمال به یک مسئله‌ی دسته‌بندی دودویی استاندارد میان داده‌های واقعی و یک توزیع نویز شناخته‌شده تبدیل گردد. فرض کنید داده‌های واقعی از توزیع $P_{\text{\lr{data}}}$ با چگالی $p_d(y)$ و داده‌های نویز از یک توزیع فرضی $P_n$ با چگالی مشخص $p_n(y)$ تولید شوند، و به ازای هر نمونه‌ی واقعی، $\nu$ نمونه‌ی نویز تولید گردد. یک دسته‌بند رگرسیون لجستیک می‌تواند احتمال تعلق یک نمونه‌ی مشاهده‌شده به توزیع داده‌ی واقعی را به‌صورت زیر بیان کند:
\begin{equation}\label{eq:chap2_nce_binary_prob}
    P(D=1 \mid y) = \frac{p_d(y)}{p_d(y) + \nu p_n(y)}
\end{equation}
گوتمان و هیواریانن اثبات کردند که اگر چگالی مدل نرمال‌نشده را با تابع انرژی پارامتربندی کنیم، کمینه‌سازی زیان آنتروپی متقاطع دودویی در این دسته‌بندی، پارامترهای مدل را به گونه‌ای همگرا می‌کند که انرژی یادگرفته‌شده متناسب با لگاریتم چگالی واقعی داده‌ها گردد، بدون آنکه نیازی به محاسبه‌ی مستقیم تابع افراز باشد.

با این حال، در مسائلی که ساختار شرطی قوی دارند (نظیر ارزیابی تناسب یک پاسخ با پرامپت)، نسخه‌ی چندکلاسه‌ی این رویکرد تحت عنوان \lr{InfoNCE}\enfootnote{information noise-contrastive estimation} که توسط ون دن اورد و همکاران (۲۰۱۸)~\cite{oord2018representation} در چارچوب یادگیری بازنمایی پیش‌بینانه معرفی شد، کارایی بسیار بالاتری نشان می‌دهد. در حالی که \lr{NCE} سنتی یک نمونه را صرفاً در برابر یک نمونه‌ی نویز مقایسه می‌کرد، \lr{InfoNCE} مدل را وادار می‌سازد تا نمونه‌ی مثبت را از میان مجموعه‌ای متشکل از یک نمونه‌ی مثبت و $K$ نمونه‌ی منفی متمایز کند.

در فرمول‌بندی مبتنی بر انرژی \lr{InfoNCE}، فرض کنید برای یک پرامپت ورودی $x_i$ و پاسخ واقعی متناظر $y_i$ (یعنی یک زوج مثبت)، مجموعه‌ای متشکل از $K$ پاسخ نامنطبق یا منفی مرتبط به‌صورت $\mathcal{N}_i = \{y'_{i, 1}, y'_{i, 2}, \dots, y'_{i, K}\}$ در اختیار باشد. با تعریف تابع امتیازدهی مبتنی بر منفیِ انرژی به‌صورت $h_\theta(x, y) = \exp(-\mathcal{E}_\theta(x, y) / \tau)$، احتمال سافت‌مکسِ نسبت‌داده‌شده به نمونه‌ی مثبت در میان این $K+1$ گزینه به‌صورت زیر تعریف می‌شود:
\begin{equation}\label{eq:chap2_infonce_prob_def}
    p(y_i \mid x_i, \mathcal{N}_i) = \frac{\exp\big(-\mathcal{E}_\theta(x_i, y_i)/\tau\big)}{\exp\big(-\mathcal{E}_\theta(x_i, y_i)/\tau\big) + \sum_{k=1}^K \exp\big(-\mathcal{E}_\theta(x_i, y'_{i, k})/\tau\big)}
\end{equation}
در این رابطه، صورت کسر امتیاز هم‌ترازیِ زوج مثبت و مخرج کسر حاصل‌جمع امتیاز هم‌ترازی نمونه‌ی مثبت و تمامی $K$ نمونه‌ی منفی است. تابع زیان \lr{InfoNCE} به‌صورت منفیِ لگاریتم این احتمال تعریف می‌گردد:
\begin{equation}\label{eq:chap2_infonce}
    \mathcal{L}_{\text{\lr{InfoNCE}}}(\theta) = -\mathbb{E} \left[ \log \frac{\exp\big(-\mathcal{E}_\theta(x_i, y_i)/\tau\big)}{\exp\big(-\mathcal{E}_\theta(x_i, y_i)/\tau\big) + \sum_{k=1}^K \exp\big(-\mathcal{E}_\theta(x_i, y'_{i, k})/\tau\big)} \right]
\end{equation}
از منظر نظریه‌ی اطلاعات، ون دن اورد و همکاران اثبات کردند که کمینه‌سازی زیان \lr{InfoNCE}، معادل با بیشینه‌سازی کران پایین اطلاعات متقابل\enfootnote{mutual information} میان متغیرهای $X$ و $Y$ است:
\begin{equation}\label{eq:chap2_mutual_info_bound}
    I(X; Y) \ge \log(K + 1) - \mathcal{L}_{\text{\lr{InfoNCE}}}
\end{equation}
در این رابطه، $I(X; Y)$ اطلاعات متقابل شانون میان دو متغیر، و $K$ تعداد نمونه‌های منفی است. این نامساوی به صراحت نشان می‌دهد که با افزایش تعداد نمونه‌های منفی ($K$)، کران پایین اطلاعات متقابل فشرده‌تر و مستحکم‌تر می‌شود و تابع انرژی یادگرفته‌شده، وابستگی‌های پنهان داده‌ها را با غنای بیشتری بازنمایی می‌کند.

بینش تحلیلی پیرامون نحوه‌ی اثرگذاری این تابع زیان، با محاسبه‌ی گرادیان آن نسبت به پارامترهای مدل ($\theta$) به‌دست می‌آید. با مشتق‌گیری زنجیره‌ای از معادله‌ی~\ref{eq:chap2_infonce}، گرادیان دقیق برای یک نمونه به فرم زیر حاصل می‌شود:
\begin{equation}\label{eq:chap2_infonce_gradient}
    \nabla_\theta \mathcal{L}_{\text{\lr{InfoNCE}}} = \frac{1}{\tau} \left[ \Big(1 - p(y_i \mid x_i, \mathcal{N}_i)\Big) \nabla_\theta \mathcal{E}_\theta(x_i, y_i) - \sum_{k=1}^K p(y'_{i, k} \mid x_i, \mathcal{N}_i) \nabla_\theta \mathcal{E}_\theta(x_i, y'_{i, k}) \right]
\end{equation}
که در آن $p(y'_{i, k} \mid x_i, \mathcal{N}_i)$ وزن احتمالاتی نرمال‌شده‌ای است که در مخرج سافت‌مکس به نمونه‌ی منفی $k$-ام اختصاص یافته است:
\begin{equation}\label{eq:chap2_infonce_neg_prob}
    p(y'_{i, k} \mid x_i, \mathcal{N}_i) = \frac{\exp\big(-\mathcal{E}_\theta(x_i, y'_{i, k})/\tau\big)}{\exp\big(-\mathcal{E}_\theta(x_i, y_i)/\tau\big) + \sum_{j=1}^K \exp\big(-\mathcal{E}_\theta(x_i, y'_{i, j})/\tau\big)}
\end{equation}
مقایسه‌ی معادله‌ی~\ref{eq:chap2_infonce_gradient} با گرادیان غیرقابل‌حل \lr{NLL} در معادله‌ی~\ref{eq:chap2_mle_gradient}، چرایی کارایی یادگیری تقابلی را به‌روشنی تبیین می‌نماید:
\begin{itemize}
    \item در فاز مثبت، گرادیان $\nabla_\theta \mathcal{E}_\theta(x_i, y_i)$ با ضریب مثبت $(1 - p(y_i \mid x_i, \mathcal{N}_i))$ ظاهر شده است که تحت گام نزول گرادیان، سطح انرژی نمونه‌ی واقعی را کاهش می‌دهد.
    \item در فاز منفی، انتگرال‌گیری ناممکن بر روی کل فضای خروجی، با یک میانگین وزنی متناهی بر روی $K$ نمونه‌ی منفی مشخص جایگزین شده است. علامت منفی پشت مجموع باعث می‌شود که گام نزول گرادیان، سطح انرژی نمونه‌های منفی را به سمت بالا هدایت کند.
    \item بر اساس خاصیت جمع احتمالات در تابع سافت‌مکس، رابطه‌ی زیر همواره برقرار است:
    \begin{equation}\label{eq:chap2_prob_balance}
        1 - p(y_i \mid x_i, \mathcal{N}_i) = \sum_{k=1}^K p(y'_{i, k} \mid x_i, \mathcal{N}_i)
    \end{equation}
    این رابطه اثبات می‌کند که مجموع نیروی دافعه‌ی اعمال‌شده بر روی نمونه‌های منفی، دقیقاً با نیروی جاذبه‌ی اعمال‌شده بر روی نمونه‌ی مثبت در تعادل است. علاوه بر این، سیستم دارای ویژگی خودتنظیمی است؛ بدین‌معنا که اگر مدل در تشخیص نمونه‌ی مثبت به اطمینان بالا برسد ($p(y_i \mid \cdot) \to 1$)، ضریب $(1 - p_0)$ و تمامی ضرایب $p_k$ به صفر میل کرده و گرادیان محو می‌شود.
\end{itemize}

این ساختار گرادیانی، اهمیت تعیین‌کننده‌ی نمونه‌های منفی سخت\enfootnote{hard negative samples} را در پایداری آموزش آشکار می‌سازد~\cite{robinson2021contrastive}. رابینسون و همکاران (۲۰۲۱) نشان دادند که اگر نمونه‌های منفی صرفاً نویزهای تصادفی یا گزینه‌های کاملاً بی‌ارتباط باشند، مدل به‌سادگی انرژی بسیار بالایی به آن‌ها نسبت می‌دهد؛ در نتیجه، مقدار $p(y'_{i, k} \mid \cdot)$ برای تمامی این نمونه‌ها به‌سرعت به صفر میل کرده و سهم آن‌ها در گرادیان معادله‌ی~\ref{eq:chap2_infonce_gradient} ناپدید می‌شود. این امر به توقف زودهنگام آموزش می‌انجامد. در مقابل، طراحی نمونه‌های منفی دشوار و شبیه به داده‌های واقعی که دارای تفاوت‌های جزئی در بافتار هستند، مقادیر $p(y'_{i, k} \mid \cdot)$ را غیرصفر نگه داشته و سیگنال‌های گرادیانی قوی و مؤثری را برای تصحیح هندسه‌ی انرژی تولید می‌نماید.



\section{انتخاب ویژگی در سطح نمونه}\label{sec:2-5}

در مبانی یادگیری ماشین، فرایند ردیابی و استخراج مهم‌ترین متغیرهای ورودی که مسئول تصمیم مدل هستند، در چارچوب مسئله‌ی انتخاب ویژگی فرمول‌بندی می‌شود. در این بخش، تفاوت بنیادین میان رویکردهای انتخاب ویژگی سراسری و انتخاب وابسته به نمونه تبیین گردیده، فرمول‌بندی ریاضی ماسک‌های خلوت ارائه شده و در نهایت، چالش مشتق‌ناپذیری و روش‌های آزادسازی پیوسته نظیر توزیع گامبل-سافت‌مکس واکاوی می‌شوند.


\subsection{تقابل انتخاب ویژگی سراسری و سطح نمونه}\label{sec:2-5-1}

به‌طور سنتی، روش‌های انتخاب ویژگی سراسری در پی یافتن یک زیرمجموعه‌ی ایستا و ثابت از ویژگی‌ها به نام $S \subset \{1, \dots, n\}$ هستند که برای تمامی نمونه‌های موجود در توزیع داده‌ها بیشترین بار اطلاعاتی را فراهم آورند (مانند اعمال جریمه‌ی $L_1$ در رگرسیون لاسو\enfootnote{LASSO} یا سنجه‌های افزایش اطلاعات در درخت‌های تصمیم). از دیدگاه نظریه‌ی اطلاعات، هدف این روش‌ها بیشینه‌سازی اطلاعات متقابل سراسری میان زیرمجموعه‌ی منتخب ورودی ($X_S$) و متغیر پاسخ ($Y$) است:
\begin{equation}\label{eq:chap2_global_fs}
    \max_{S \subset \{1, \dots, n\}, \, |S| \le k} I(X_S; Y)
\end{equation}
در این فرمول‌بندی، زیرمجموعه‌ی $S$ برای تمامی ورودی‌ها یکسان است. با این وجود، در داده‌های پیچیده و بافتارمحور نظیر متون طبیعی، میزان اهمیت یک متغیر نمی‌تواند به‌صورت سراسری ثابت باشد. برای مثال، کلمه‌ای خاص ممکن است در یک متن کلیدی‌ترین نقش تشخیصی را ایفا کند، اما در متنی دیگر با موضوعی متفاوت، کاملاً حاشیه‌ای و بی‌اثر باشد.

برای رفع این محدودیت، پارادایم «انتخاب ویژگی در سطح نمونه»\enfootnote{instance-wise feature selection} توسط چن و همکاران (۲۰۱۸) در چارچوب \lr{L2X}~\cite{chen2018learning} و یون و همکاران (۲۰۱۹) در روش \lr{INVASE}~\cite{yoon2019invase} فرمول‌بندی گردید. در این دیدگاه، انتخاب ویژگی به‌صورت یک تابع شرطی و وابسته به ورودی به نام مفسر ($\mathcal{IN}$) تعریف می‌شود که برای هر نمونه‌ی خاص $X=x$، مهم‌ترین زیرمجموعه‌ی متناظر با همان نمونه ($S(x)$) را شناسایی می‌کند. به‌بیان ریاضی، هدف بیشینه‌سازی اطلاعات متقابل شرطی تحت نگاشت مفسر است:
\begin{equation}\label{eq:chap2_local_fs}
    \max_{\mathcal{IN}} I(X_{\mathcal{IN}(X)}; Y)
\end{equation}
چن و همکاران اثبات کردند که برای یک نمونه‌ی مشاهده‌شده‌ی معین ($X=x$)، بهینه‌سازی این معیار معادل با یافتن زیرمجموعه‌ای است که واگرایی کولبک-لایبلر میان توزیع پیش‌بینی بر پایه‌ی کل ویژگی‌ها و توزیع پیش‌بینی بر اساس ویژگی‌های منتخب را کمینه کند:
\begin{equation}\label{eq:chap2_l2x_kl_opt}
    S^*(x) = \arg\min_{S: |S| \le k} D_{\text{\lr{KL}}}\big(P(Y \mid X=x) \parallel P(Y \mid X_S = x_S)\big)
\end{equation}
در این رابطه، $P(Y \mid X=x)$ توزیع پیش‌بینی مدل با دریافت تمامی ویژگی‌ها، و $P(Y \mid X_S = x_S)$ توزیع پیش‌بینی صرفاً با دسترسی به زیرمجموعه‌ی منتخب $S$ است. این فرمول‌بندی نشان می‌دهد که انتخاب ویژگی در سطح نمونه، مستقیماً به دنبال یافتن ویژگی‌هایی است که بدون حضور آن‌ها، رفتار مدل در قبال آن نمونه‌ی مشخص بیشترین انحراف را تجربه می‌کند.


\subsection{فرمول‌بندی ریاضی ماسک دودویی و قید خلوتی}\label{sec:2-5-2}

از دیدگاه محاسباتی، فرایند انتخاب ویژگی بر روی یک ورودی مشخص را می‌توان با استفاده از یک عملگر ضرب تانسوری و ماسک دودویی مدل‌سازی نمود. فرض کنید داده‌ی ورودی شامل $n$ ویژگی برداری پیوسته باشد که در قالب ماتریس $X \in \mathbb{R}^{n \times d}$ مرتب شده‌اند (به‌عنوان مثال، $n$ بردار تعبیه‌ی جمله به‌عنوان پیش‌آیندهای گزاره‌ای ورودی به طول $d$).

یک بردار ماسک اهمیت دودویی به‌صورت $\mathbf{m} = {(m_1, m_2, \dots, m_n)}^\top \in {\{0, 1\}}^n$ تعریف می‌گردد. زیرمجموعه‌ی ویژگی‌های انتخاب‌شده ($X_S$) از طریق عملگر ضرب درایه‌به‌درایه‌ی هادامارد با پخش‌پذیری در طول بعد پنهان $d$ استخراج می‌شود:
\begin{equation}\label{eq:chap2_masking}
    X_S = \mathbf{m} \odot X
\end{equation}
در این رابطه، اگر درایه‌ی $m_i = 1$ باشد، بردار ویژگی سطر $i$-ام بدون تغییر حفظ می‌شود، و اگر $m_i = 0$ باشد، تمام مؤلفه‌های آن سطر به صفر تبدیل شده و اثر آن از محاسبات بعدی حذف می‌گردد.

تعیین ماسک بهینه نیازمند ایجاد توازن میان دو هدف متضاد است: نخست، اصل کفایت؛ به‌گونه‌ای که اطلاعات موجود در داده‌ی ماسک‌شده ($X_S$) برای هدایت مدل به سوی بازتولید خروجی هدف کافی باشد و خروجی مدل روی داده‌ی ماسک‌شده ($f(X_S)$) به خروجی داده‌ی کامل ($f(X)$) نزدیک بماند؛ و دوم، اصل ضرورت و فشردگی، مبنی بر به‌حداقل رساندن تعداد ویژگی‌های انتخاب‌شده تا انتساب حاصل فشرده، شفاف و انسان‌خوانا باشد. این تعادل در قالب مسئله‌ی بهینه‌سازی مقید به خلوتی زیر بیان می‌شود:
\begin{equation}\label{eq:chap2_sparsity_opt}
    \min_{\mathbf{m} \in {\{0, 1\}}^n} \mathcal{L}\big(f(X), f(\mathbf{m} \odot X)\big) + \lambda \|\mathbf{m}\|_0
\end{equation}
در این معادله، $\mathcal{L}$ یک تابع زیان برای سنجش میزان انحراف یا خطای بازسازی خروجی مدل است. عبارت دوم شامل شبه‌نرم صفر یعنی $\|\mathbf{m}\|_0 = \sum_{i=1}^n \mathbb{I}(m_i \neq 0)$ است که تعداد درایه‌های غیرصفر بردار ماسک را شمارش می‌کند، و $\lambda > 0$ ابرپارامتر جریمه‌ی پیچیدگی است که شدت خلوت‌سازی را تنظیم می‌نماید.

استپین و همکاران (۲۰۲۱)~\cite{stepin2021survey} در بررسی جامع روش‌های تبیین پادواقعی و تفسیرپذیری جعبه‌سیاه، تبیین می‌کنند که حداقل‌سازی شبه‌نرم صفر تجلی «اصل تیغ اوکام»\enfootnote{Occam's razor} است؛ زیرا یک توضیح زمانی معتبر و ارزشمند تلقی می‌شود که فاقد هرگونه متغیر زائد یا افزونگی تصادفی باشد.


\subsection{چالش مشتق‌ناپذیری و آزادسازی گامبل-سافت‌مکس}\label{sec:2-5-3}

حل مستقیم مسئله‌ی بهینه‌سازی در معادله‌ی~\ref{eq:chap2_sparsity_opt} بر روی فضای گسسته‌ی ماسک‌های دودویی ${\{0, 1\}}^n$، یک مسئله‌ی ترکیبیاتی با فضای حالاتی به اندازه‌ی $2^n$ است که ارزیابی جامع آن برای مقادیر متعارف $n$ از نظر محاسباتی ناممکن است. از سوی دیگر، تلاش برای بهینه‌سازی این ماسک با استفاده از شبکه‌های عصبی پارامتریک و الگوریتم پس‌انتشار با مانع بنیادین مشتق‌ناپذیری روبرو می‌شود؛ زیرا مشتق عملگرهای گسسته‌ساز و توابع پله‌ای نسبت به ورودی‌های پیوسته در تمام نقاط، صفر یا تعریف‌نشده ($\nabla_{\mathbf{z}} \mathbf{m} = \mathbf{0}$) است. در نتیجه، گرادیان سیگنال خطا نمی‌تواند از متغیرهای گسسته عبور کرده و پارامترهای شبکه‌ی مفسر را به‌روزرسانی کند~\cite{chen2018learning}.

برای غلبه بر این چالش، مفهوم آزادسازی پیوسته\enfootnote{continuous relaxation} مطرح می‌گردد که متغیرهای تصادفی گسسته را با توابع هموار، پیوسته و مشتق‌پذیر تقریب می‌زند. برجسته‌ترین روش در این زمینه، ترفند گامبل-سافت‌مکس\enfootnote{Gumbel-softmax} (یا توزیع کانکریت\enfootnote{Concrete distribution}) است که به‌صورت هم‌زمان و مستقل توسط جانگ و همکاران (۲۰۱۷)~\cite{jang2017categorical} و مدیسون و همکاران (۲۰۱۷)~\cite{maddison2017concrete} معرفی گردید.

این روش بر پایه‌ی تعمیم هموارِ ترفند گامبل-مکس\enfootnote{Gumbel-max trick} بنا شده است. در ترفند گامبل-مکس کلاسیک، برای نمونه‌برداری گسسته از یک توزیع دسته‌بندی با پارامترهای لاگ-احتمال (یعنی لوجیت‌ها) به‌صورت $\mathbf{z} = (z_1, z_2, \dots, z_n) \in \mathbb{R}^n$، ابتدا متغیرهای تصادفی مستقل نویز استاندارد گامبل ($g_i$) با استفاده از متغیرهای یکنواخت پیوسته ($\xi_i \sim \text{\lr{Uniform}}(0, 1)$) به‌صورت زیر تولید می‌شوند:
\begin{equation}\label{eq:chap2_gumbel_noise}
    g_i = -\log\big(-\log(\xi_i)\big) \quad \forall i \in \{1, \dots, n\}
\end{equation}
سپس، اندیس متغیر منتخب از طریق عملگر $\arg\max_i (z_i + g_i)$ استخراج می‌گردد. متغیر تصادفی حاصل دقیقاً دارای توزیع دسته‌بندی با احتمالات متناسب با $\exp(z_i)$ است، اما عملگر $\arg\max$ همچنان مشتق‌ناپذیر است.

نوآوری جانگ و همکاران و مدیسون و همکاران، جایگزینی عملگر ناپیوسته‌ی $\arg\max$ با عملگر هموار و پیوسته‌ی سافت‌مکس به‌همراه یک پارامتر دما (به‌صورت $\tau > 0$) بود. بدین‌ترتیب، بردار پیوسته‌ی $\mathbf{c} = {(c_1, c_2, \dots, c_n)}^\top$ بر روی سیمپلکس احتمالاتی با ضابطه‌ی زیر تعریف می‌گردد:
\begin{equation}\label{eq:chap2_gumbel_softmax}
    c_i = \frac{\exp\big((z_i + g_i) / \tau\big)}{\sum_{j=1}^n \exp\big((z_j + g_j) / \tau\big)} \quad \forall i \in \{1, \dots, n\}
\end{equation}
در این معادله، $z_i$ لوجیت تولیدشده توسط شبکه‌ی عصبی، $g_i$ نویز تصادفی گامبل، و $\tau$ پارامتر دما است. رفتار این تابع با تغییر دما به‌صورت زیر تنظیم می‌شود:
\begin{itemize}
    \item هنگامی که دما به سمت صفر میل می‌کند ($\tau \to 0$)، خروجی $\mathbf{c}$ به سمت یک بردار یک‌گانه‌ی قطعی و متعامد میل می‌کند که درایه‌ی متناظر با بیشینه‌ی $(z_i + g_i)$ در آن برابر یک و سایر درایه‌ها صفر خواهند بود.
    \vspace{-0.15cm}
    \item هنگامی که دما به سمت بی‌نهایت میل کند ($\tau \to \infty$)، خروجی به سمت یک بردار یکنواخت با مقادیر یکسان ($1/n$) میل می‌نماید.
    \vspace{-0.15cm}
    \item برای مقادیر متداول $\tau > 0$، خروجی $\mathbf{c}$ در فضای درون سیمپلکس $\Delta^{n-1}$ قرار گرفته، نسبت به تمامی لوجیت‌های $\mathbf{z}$ مشتق‌پذیر پیوسته است، و به گرادیان‌ها اجازه می‌دهد تا با استفاده از ترفند بازپارامتربندی\enfootnote{reparameterization trick} به‌صورت مستقیم در طول گراف محاسباتی منتشر شوند.
\end{itemize}

چن و همکاران (۲۰۱۸) در روش \lr{L2X}~\cite{chen2018learning} از این ساختار آزادسازی (و تعمیم آن به نمونه‌برداری از زیرمجموعه‌های $k$-تایی) بهره گرفتند تا فرایند انتخاب ویژگی در سطح نمونه را به یک مسئله‌ی بهینه‌سازی انتها‌به‌انتها و قابل‌آموزش با روش‌های نزول گرادیان تصادفی تبدیل کنند.
